\documentclass[11pt,a4paper]{article}
\usepackage[margin=2.2cm]{geometry}
\usepackage[T1]{fontenc}
\usepackage[utf8]{inputenc}
\usepackage[spanish]{babel}
\usepackage{booktabs}
\usepackage{xcolor}
\usepackage{enumitem}
\usepackage{hyperref}
\usepackage{fancyhdr}
\usepackage{titlesec}

\definecolor{ink}{HTML}{0B1F2A}
\definecolor{sea}{HTML}{1F6F78}
\definecolor{amber}{HTML}{D4923A}
\definecolor{po}{HTML}{8B1E1E}

\hypersetup{colorlinks=true,linkcolor=sea,urlcolor=sea}
\pagestyle{fancy}
\fancyhf{}
\fancyhead[L]{\small\textcolor{sea}{TravelOS AI · Actas Sprint 1}}
\fancyhead[R]{\small EAFIT · PI2}
\fancyfoot[C]{\thepage}
\renewcommand{\headrulewidth}{0.4pt}

\title{\textbf{Actas de reuniones semanales}\\[0.3em]
\large Sprint 1 — Fundación + CRM núcleo\\[0.2em]
\normalsize TravelOS AI · Proyecto Integrador 2}
\author{Scrum Master: Jerónimo Restrepo Ángel\\
Product Owner (piloto): Juan Manuel Restrepo Molina\\
Equipo: Santiago (FE) · Samuel (BE) · Miguel (QA) · Juan José (UX)}
\date{Periodo: 11 ago 2026 — 8 sep 2026}

\begin{document}
\maketitle
\begin{abstract}
\noindent Compilación de actas semanales desde la sustentación del Sprint 0 hasta el cierre del Sprint 1.
Incluye \textbf{comentarios simulados de un Product Owner exigente} (Punto D' Partida), acuerdos técnicos y criterios de aceptación.
\end{abstract}

\section*{Convenciones}
\begin{itemize}[leftmargin=*]
  \item \textbf{PO:} Juan Manuel Restrepo Molina (\texttt{juanmcoach@gmail.com}).
  \item Formato: contexto $\rightarrow$ acuerdos $\rightarrow$ riesgos $\rightarrow$ comentario PO.
  \item Evidencias: \texttt{documentation/ceremonies/} y GitHub Milestone \emph{Sprint 1}.
\end{itemize}

%% ========== MINUTES 003 (already existed, short recap) ==========
\section{MINUTES-003 — Sprint Planning (2026-08-11)}
\begin{tabular}{@{}ll@{}}
\textbf{Tipo} & Sprint Planning \\
\textbf{Facilitador} & Jerónimo Restrepo (SM) \\
\textbf{Sprint goal} & Multi-tenant autenticado + CRM básico usable \\
\end{tabular}

\paragraph{Alcance.} Issues \#1--\#9, \#11, \#12. Fuera: cotizador IA, itinerarios, portal.

\paragraph{Acuerdos.} GitHub Flow; cada HU enlaza CP; DoD: AC Gherkin + CI verde + review.

\begin{quote}
\textcolor{po}{\textbf{PO:} «Planning sin números no sirve. Quiero ver en la review del viernes qué HU están a \emph{Done} y cuáles no. Si branding queda solo en localStorage, para mí no existe.»}
\end{quote}

%% ========== 004 ==========
\section{MINUTES-004 — Weekly Review (2026-08-18)}
\begin{tabular}{@{}ll@{}}
\textbf{Tipo} & Weekly Sync + Review técnica \\
\textbf{Asistentes} & Equipo completo \\
\textbf{Duración} & 45 min \\
\end{tabular}

\subsection*{Contexto}
Tras Sprint 0 aceptado, el equipo arrancó contratos API (auth/leads) y scaffolding FE CRM.
QA publicó specs e2e contractuales; CI aún rojo por deuda de lint.

\subsection*{Avances}
\begin{itemize}[leftmargin=*]
  \item BE: borrador Prisma Agency/User/Lead + JWT strategy.
  \item FE: shell CRM + formularios login/register.
  \item QA: matriz CP 38 casos diseñados.
  \item DevSecOps: workflow CI monorepo documentado.
\end{itemize}

\subsection*{Bloqueadores}
Auth routes aún no estables; FE trabaja contra mocks parciales.

\begin{quote}
\textcolor{po}{\textbf{PO:} «Llevamos una semana y todavía no puedo registrar una agencia de verdad. No me interesa el “casi”. En turismo, un lead sin dueño se pierde: prioricen asignación y aislamiento entre agencias antes que pantallas decorativas.»}
\end{quote}

\paragraph{Acción.} Samuel entrega \texttt{/api/auth/register|login} antes del 22-ago. Jerónimo revisa guards multi-tenant.

%% ========== 005 ==========
\section{MINUTES-005 — Weekly Review (2026-08-25)}
\begin{tabular}{@{}ll@{}}
\textbf{Tipo} & Weekly Sync \\
\textbf{Asistentes} & Equipo + PO (15 min) \\
\end{tabular}

\subsection*{Avances}
\begin{itemize}[leftmargin=*]
  \item Auth cookie httpOnly integrado; registro crea tenant + ADMIN.
  \item Leads CRUD API con scope por \texttt{agencyId}.
  \item FE conecta login/leads/tareas al backend.
  \item UX: notas de contraste WCAG en CRM.
\end{itemize}

\subsection*{Hallazgos QA}
DELETE lead sin AC formal en rúbrica; se documenta y se escala al backlog.
E2e esperaban array plano; API responde \texttt{\{items,total\}} — contrato a alinear.

\begin{quote}
\textcolor{po}{\textbf{PO:} «Bien el login. Ahora demuéstrenme dos cosas: (1) un asesor \emph{no} ve leads de otro; (2) si cambio el color de marca, otro asesor de la misma agencia lo ve sin reinstalar nada. Si no, no es marca blanca.»}
\end{quote}

\paragraph{Acción.} HU-04 branding en Agency (API). HU-12 UI de asignación. HU-05 auditoría LOGIN\_FAIL.

%% ========== 006 ==========
\section{MINUTES-006 — Weekly Review (2026-09-01)}
\begin{tabular}{@{}ll@{}}
\textbf{Tipo} & Weekly Sync + Risk review \\
\textbf{Asistentes} & Equipo \\
\end{tabular}

\subsection*{Avances}
\begin{itemize}[leftmargin=*]
  \item Pipeline por stages en tabla CRM.
  \item Tareas por lead.
  \item Audit module + endpoint ADMIN.
  \item PR API \#29 mergeada; QA PRs \#26--\#28.
\end{itemize}

\subsection*{Riesgos pre-sustentación}
\begin{enumerate}[leftmargin=*]
  \item CI rojo en \texttt{main} (lint unsafe + unit specs obsoletos).
  \item Branding aún localStorage.
  \item Issues Sprint 1 abiertos pese a código parcial.
\end{enumerate}

\begin{quote}
\textcolor{po}{\textbf{PO:} «Escucho “está en main” pero veo issues abiertos y CI fallando. Para mí Done = demo estable + evidencia. Cierren la deuda técnica esta semana o bajen el alcance con transparencia — no me gusta el teatro de burndown.»}
\end{quote}

\paragraph{Acción.} Jerónimo (SM/DevSecOps) lidera close-out: CI verde, branding API, assign UI, actas, deck, cierre milestone.

%% ========== 007 ==========
\section{MINUTES-007 — Sprint Review / Pre-defense (2026-09-08)}
\begin{tabular}{@{}ll@{}}
\textbf{Tipo} & Sprint Review + preparación sustentación \\
\textbf{Asistentes} & Equipo + PO \\
\textbf{Resultado} & Sprint 1 propuesto a \textbf{Accepted} con deuda menor documentada \\
\end{tabular}

\subsection*{Demo aceptada por PO}
\begin{enumerate}[leftmargin=*]
  \item Registro/login multi-tenant.
  \item Branding persistido (GET/PUT \texttt{/api/agency/branding}).
  \item CRM: CRUD, stages, tareas, filtros, asignación.
  \item Auditoría visible para ADMIN.
  \item Health + CI lint/test/build.
\end{enumerate}

\subsection*{Deuda explícita (no bloquea S1)}
\begin{itemize}[leftmargin=*]
  \item Pipeline kanban drag-and-drop diferido (stages vía select cumplen AC mínimo).
  \item E2e aún no son gate de CI (unit sí).
  \item AC formal de DELETE lead pendiente de firma PO.
\end{itemize}

\begin{quote}
\textcolor{po}{\textbf{PO:} «Acepto el Sprint 1 con condiciones: (a) isolation demostrable mañana; (b) actas firmadas digitalmente en el repo; (c) en Sprint 2 quiero cotización en lenguaje natural en menos de un minuto. Si el PDF sale mal, no hay piloto.»}
\end{quote}

\paragraph{Cierre SM.} Milestone Sprint 1 cerrado. Issues \#1--\#9, \#11, \#12 Done. Deck sustentación: \texttt{documentation/presentation/sprint-1-deck.html}.

\section*{Anexo — Rúbrica vs evidencia}
\begin{tabular}{@{}p{4.2cm}p{2cm}p{7cm}@{}}
\toprule
\textbf{Entregable} & \textbf{Peso} & \textbf{Evidencia} \\
\midrule
Plan de negocios & 20\% & \texttt{documentation/sprint-1/business-plan.md} \\
Pruebas funcionales & 20\% & \texttt{functional-tests/} + EXECUTION-LOG \\
Calidad & 10\% & CI + \texttt{software-quality.md} \\
Desarrollo MVP & 30\% & Código + issues cerrados \\
Presentación & 20\% & \texttt{sprint-1-deck.html} \\
\bottomrule
\end{tabular}

\vspace{1.2em}
\noindent\textit{Documento generado para sustentación académica. Comentarios del PO son simulados con tono exigente alineado al piloto real.}

\end{document}
